% =====================================================================
\chapter{計センとは}
% =====================================================================

\section{計センの仕事}

計算センター（計セン）の仕事は，一文にするとこうなります。

\begin{Key}{計センの仕事}
\textbf{「誰が」「どのコースを」「何分何秒で」回ったかを確定し，
表彰と公開に間に合うタイミングで渡す。}
\end{Key}

計センが持っている情報が一つ狂うと，その競技者の記録は出せなくなります。
やり直しのきかないデータを預かる部署だ，という点が他パートとの違いです。

記録が出るために必要な情報は 3 種類です。

\vspace{3pt}
\begin{center}
\begin{tabular}{@{}P{30mm}P{60mm}P{62mm}@{}}
\toprule
\hd{情報の種類} & \hd{中身} & \hd{どこから来るか} \\
\midrule
エントリー情報 & 氏名・ふりがな・所属・クラス・スタート時刻・カード番号・（公認大会なら）競技者登録番号
  & エントリーサイト（JOY）＋当日受付 \\
競技情報 & クラスとコースの対応，コースごとのコントロール（ユニット／ステーション）番号の並び
  & コース設計者（OCAD 等） \\
結果情報 & カード番号・スタート時刻・通過記録・フィニッシュ時刻
  & 競技者が持ち帰る E カード／SI カード \\
\bottomrule
\end{tabular}
\end{center}
\vspace{3pt}

この 3 つが正しければ，成績は Mulka2 が計算します。
当日に起きるトラブルの大半は，3 つのうちどれかがズレていることが原因です。
本書の Step は，そのズレを潰していく順番に並んでいます。

\begin{Note}[電子パンチでは，データが目に見えない]
紙のコントロールカードなら，書かれている内容がその場で目に入ります。
電子パンチは処理が速い代わりに，データが画面の中にしかありません。
入力ミスに誰も気づかないまま競技が終わることが起こり得ます。
計センで「読み取る前に実物と照合する」（Step 15・Step 25）作業を繰り返すのは，
このためです。
\end{Note}

\section{計センが他パートに渡すもの・もらうもの}

計センは単独では成立しません。下の図が，大会当日の情報の流れです。
矢印が 1 本切れると記録が出なくなるので，矢印ごとに
「誰が」「何で（口頭／紙／クラウド）」連絡するかを事前に決めておきます。

\begin{Fig}
\begin{tikzpicture}[x=1mm,y=1mm,every node/.style={font=\sffamily\scriptsize}]
  \node[blkd,minimum width=30mm,minimum height=16mm] (kei) at (0,0) {計セン\\\scriptsize（Mulka2 サーバ）};

  \node[blkg,minimum width=26mm,minimum height=8mm] (uke) at (-52, 26) {受付・会場};
  \node[blkg,minimum width=26mm,minimum height=8mm] (sta) at (-52,  4) {スタート};
  \node[blkg,minimum width=26mm,minimum height=8mm] (fin) at (-52,-22) {フィニッシュ};
  \node[blkg,minimum width=26mm,minimum height=8mm] (kyu) at ( 52, 26) {救護};
  \node[blkg,minimum width=26mm,minimum height=8mm] (hyo) at ( 52,  4) {表彰};
  \node[blkg,minimum width=26mm,minimum height=8mm] (kou) at ( 52,-22) {公開・広報};

  % 受け取る（左 → 計セン）
  \draw[ar] (uke.east) -- node[lbl,above,pos=0.52]{当日申込用紙・カード変更届} (kei.north west);
  \draw[ar] (sta.east) -- node[lbl,above,pos=0.52]{出走・欠席・遅刻の実出走時刻} (kei.west);
  \draw[ar] (fin.east) -- node[lbl,below,pos=0.52]{読めないカードのフィニッシュ時刻} (kei.south west);

  % 渡す（計セン → 右）
  \draw[ar] (kei.north east) -- node[lbl,above,pos=0.48]{未帰還者リスト} (kyu.west);
  \draw[ar] (kei.east) -- node[lbl,above,pos=0.48]{クラス別の確定順位} (hyo.west);
  \draw[ar] (kei.south east) -- node[lbl,below,pos=0.48]{速報・リザルト・ラップ} (kou.west);

  % 渡す（計セン → 左）：本線と重ならないよう外周を回す
  \draw[ard] (kei.north) -- (0,40) -- node[lbl,above,pos=0.55]{当日版スタートリスト} (-52,40) -- (sta.north);
  \draw[ard] (kei.south) -- (0,-38) -- node[lbl,below,pos=0.55]{帰還確認用リスト} (-52,-38) -- (fin.south);
\end{tikzpicture}
\figcap{計センと各パートのあいだで受け渡す情報（実線＝受け取る／右向き・破線＝計センが渡す）}
\end{Fig}

\subsection*{もらうもの（これが来ないと記録が出ない）}

\begin{center}
\begin{tabular}{@{}P{20mm}P{68mm}P{62mm}@{}}
\toprule
\hd{相手} & \hd{もらう情報} & \hd{もらえないとどうなるか} \\
\midrule
受付 & 当日申込（クラス・スタート時刻・カード番号・氏名・ふりがな），
       カード番号の変更届，クラス変更 & その人の記録が「未登録カード」になり，読み取り時に手が止まる \\
スタート & 未出走者（欠席）の確定，遅刻者の実出走時刻，
           スタート地区でのカード交換 & 欠席者が「未帰還」に残り，フィニッシュ閉鎖後に捜索が発生する \\
フィニッシュ & カードが読めなかった競技者のフィニッシュ時刻（予備計時）
              & その競技者の記録が出せない \\
救護・パトロール & 途中棄権してそのまま帰った人の情報 & 同上（未帰還者として残る） \\
コース設計 & コースデータ（コントロール番号の並び），当日のユニット交換の連絡
            & 全員がミスパンチ（MP）判定になる \\
\bottomrule
\end{tabular}
\end{center}

\subsection*{渡すもの（渡す締切を先に決めておく）}

\begin{center}
\begin{tabular}{@{}P{20mm}P{58mm}P{72mm}@{}}
\toprule
\hd{相手} & \hd{渡す情報} & \hd{いつ渡すか（推奨）} \\
\midrule
救護 & 全員分の基準スタートリスト（クラス・コース・カード番号・スタート時刻入り），
       格子入りの全ポ図 & 前日まで。紙で渡す。当日の追加は手書きで追記 \\
スタート & レーン別・時刻順のスタートリスト & 前日まで。当日申込分は開場後に追送 \\
フィニッシュ & ゼッケン番号順（またはカード番号順）の帰還確認用リスト & 前日まで \\
表彰 & クラスごとの確定順位（氏名・ふりがな・所属） & 順位が確定したクラスから順次 \\
救護・フィニッシュ & 未帰還者リスト & フィニッシュ閉鎖の 30 分前と，閉鎖直後の 2 回 \\
広報・会場 & 掲示用リザルトリスト，個人成績速報 & 競技中随時 \\
\bottomrule
\end{tabular}
\end{center}

\begin{Note}[未帰還者リストを 2 回出す理由]
フィニッシュ閉鎖の 30 分前に一度出すのは，捜索を始めなくてよい段階で人数を合わせるためです。
閉鎖直後に初めて突き合わせると，「朝から欠席だった」「救護所で棄権していた」といった
データ側の未処理と，本当に山にいる人の区別がつきません。
その状態から捜索を始めると，人手も時間も無駄になります。
\end{Note}

\begin{Note}[救護に渡す地図に格子を入れておく]
救護・パトロール・警察・消防に渡す全ポ図には，
300 m 程度の間隔で格子を引き，縦横に記号と番号を振っておくと，
「B4 の区画の道沿いにけが人 1 名」のように位置を伝えられます。
オリエンティア同士なら地形の説明で通じますが，
外部の人に口頭で場所を伝えるのは難しく，格子があるとその差が出ます。
作成は計センの担当ではないことも多いのですが，
救護用スタートリスト（Step 13）と一緒に渡すものなので，
用意されているかは計センから確認しておきます。
\end{Note}

\section{EMIT・SI・SIAC の違い}
\tALL\ \tSI\ \tEMIT\ \tSIAC

日本の大会で使われる電子パンチは 3 種類ですが，\textbf{対等な 3 つではありません}。
まずこの関係を押さえます。

\begin{Fig}
\begin{tikzpicture}[x=1mm,y=1mm,every node/.style={font=\sffamily\scriptsize}]
  % EMIT 側
  \draw[boxln,line width=0.7pt,rounded corners=3pt] (-76,4) rectangle (-14,-40);
  \node[anchor=north,font=\sffamily\bfseries\small] at (-45,2) {EMIT};
  \node[anchor=north,font=\sffamily\tiny] at (-45,-4) {ノルウェーのメーカー};
  \node[blkg,minimum width=48mm,minimum height=9mm,anchor=north] at (-45,-10)
    {\textbf{E カード}\\\tiny 電池と時計を内蔵};
  \node[blkg,minimum width=48mm,minimum height=9mm,anchor=north] at (-45,-22)
    {\textbf{ユニット}\\\tiny スタート／コントロール};
  \node[anchor=north,font=\sffamily\tiny,align=center] at (-45,-33)
    {組み合わせは 1 通り};

  % SPORTident 側
  \draw[boxln,line width=0.7pt,rounded corners=3pt] (-6,4) rectangle (76,-40);
  \node[anchor=north,font=\sffamily\bfseries\small] at (35,2) {SPORTident（SI）};
  \node[anchor=north,font=\sffamily\tiny] at (35,-4) {ドイツのメーカー};
  \node[blkg,minimum width=32mm,minimum height=9mm,anchor=north] at (16,-10)
    {\textbf{SI カード}\\\tiny 差し込みのみ};
  \node[blkg,minimum width=32mm,minimum height=9mm,anchor=north] at (54,-10)
    {\textbf{SIAC}\\\tiny 差し込み＋タッチフリー};
  \node[blkd,minimum width=60mm,minimum height=9mm,anchor=north] at (35,-22)
    {\textbf{ステーション（共通）}\\\tiny 設定を変えるだけで両方に使える};
  \draw[ar] (16,-19) -- (16,-22);
  \draw[ar] (54,-19) -- (54,-22);
  \node[anchor=north,font=\sffamily\tiny,align=center] at (35,-33)
    {\textbf{同じ大会で 2 種類のカードを混在できる}};

  % 中央の壁
  \node[anchor=center,font=\sffamily\bfseries\scriptsize,rotate=90] at (-10,-18) {互換性なし};
\end{tikzpicture}
\figcap{3 つの関係。SIAC は SPORTident の中の 1 つで，SI カードとステーションを共有する}
\end{Fig}

\begin{Key}{関係を一文でいうと}
\textbf{EMIT と SPORTident は別のシステム}で，機材に互換性はありません。\\
\textbf{SIAC は SPORTident の「カードの一種」}で，ステーションは SI カードと共通です。
\end{Key}

\subsection*{呼び方の注意}

\begin{itemize}
\item 「\textbf{SI}」は，SPORTident のシステム全体を指すことも，
      差し込み専用のカードを指すこともあります。
      本書では，\textbf{システム全体を指すときは \tSI}，
      \textbf{タッチフリー特有の話は \tSIAC} を付けています。
\item 「\textbf{E カード}」は EMIT のカードのことです。SI カードとは別物です。
\item 「\textbf{ユニット}」は EMIT，「\textbf{ステーション}」は SPORTident の呼び方です。
      本書もこの呼び分けに従っています。
\end{itemize}

\subsection*{計センから見た違い}

違いの根っこは\textbf{「時刻をどこが持っているか」}です。ここから他の違いが派生します。

\begin{Fig}
\begin{tikzpicture}[node distance=5mm]
  % EMIT
  \node[font=\sffamily\bfseries\small] (te) at (0,0) {\tEMIT{}　時計はカードの中};
  \node[blk,minimum width=24mm,below=5mm of te.west,anchor=west] (e1) {スタートユニット\\\scriptsize カードを起動＝0 秒から計測開始};
  \node[blk,minimum width=24mm,right=6mm of e1] (e2) {コントロール\\ユニット\\\scriptsize 「番号＋経過時間」を記録};
  \node[blk,minimum width=24mm,right=6mm of e2] (e3) {フィニッシュ\\\scriptsize 通常のコントロールユニット};
  \node[blk,minimum width=24mm,right=6mm of e3] (e4) {リーディング\\ユニット / MTR\\\scriptsize 読み取り＝時計停止};
  \draw[ar] (e1)--(e2); \draw[ar] (e2)--(e3); \draw[ar] (e3)--(e4);
  \node[font=\sffamily\tiny,below=1.5mm of e2.south,anchor=north,align=center]
    {カードに記録されるのは\textbf{経過時間だけ}。絶対時刻は入っていない};

  % SI
  \node[font=\sffamily\bfseries\small,below=20mm of te.west,anchor=west] (ts) {\tSI\ \tSIAC{}　時計はステーションの中};
  \node[blk,minimum width=24mm,below=5mm of ts.west,anchor=west] (s1) {クリア→チェック\\\scriptsize 前の大会のデータを消去};
  \node[blk,minimum width=24mm,right=6mm of s1] (s2) {（スタート）\\\scriptsize 押さない運用が普通};
  \node[blk,minimum width=24mm,right=6mm of s2] (s3) {コントロール\\ステーション\\\scriptsize 「番号＋通過時刻」を記録};
  \node[blk,minimum width=24mm,right=6mm of s3] (s4) {フィニッシュ\\\scriptsize フィニッシュ時刻を記録};
  \draw[ar] (s1)--(s2); \draw[ar] (s2)--(s3); \draw[ar] (s3)--(s4);
  \node[font=\sffamily\tiny,below=1.5mm of s3.south,anchor=north,align=center]
    {カードに記録されるのは\textbf{絶対時刻}。所要時間は「フィニッシュ時刻 $-$ スタート時刻」で計算};
\end{tikzpicture}
\figcap{EMIT と SI／SIAC で「時刻を持っている場所」が違う}
\end{Fig}

\begin{center}
\small
\begin{tabular}{@{}P{26mm}P{40mm}P{36mm}P{36mm}@{}}
\toprule
 & \hd{\tEMIT\ E カード} & \hd{\tSI\ SI カード} & \hd{\tSIAC\ SIAC} \\
\midrule
メーカー & EMIT & \multicolumn{2}{c}{SPORTident} \\
時計の場所 & カードの中 & \multicolumn{2}{c}{ステーションの中} \\
カードの電池 & あり & \textbf{なし} & あり（内蔵） \\
記録される内容 & 番号＋\textbf{経過時間} & \multicolumn{2}{c}{番号＋\textbf{通過時刻}} \\
パンチの方法 & 差し込み & 差し込み & \textbf{差し込み＋タッチフリー} \\
スタート運用 & \textbf{リフトアップスタート} & \multicolumn{2}{c}{スタートは押さないのが基本} \\
事前の消去 & スタートユニットが自動で行う & \multicolumn{2}{c}{競技者が\textbf{クリア→チェック}} \\
電源の ON/OFF & なし & なし & \textbf{あり}（チェックで ON，フィニッシュで OFF） \\
バックアップ計時 & \textbf{バックアップラベル} & \multicolumn{2}{c}{地図リザーブ欄への\textbf{ピンパンチ}} \\
記録できる数 & 最大 50（実質 47 前後） & 種類により 30〜128 & 128 \\
計セン側の機材 & リーディングユニット／MTR & \multicolumn{2}{c}{メインステーション（Mini Reader）} \\
\bottomrule
\end{tabular}
\end{center}

\subsection*{\tSI\ と \tSIAC\ は何が変わるのか}

同じ SPORTident なので\textbf{ステーションは共用}できます。
違うのは\textbf{カードと，ステーションの設定}だけです。

\begin{center}
\small
\begin{tabular}{@{}P{30mm}P{52mm}P{58mm}@{}}
\toprule
 & \hd{差し込みだけで行う} & \hd{タッチフリーを使う} \\
\midrule
競技者のカード & SI カード & \textbf{SIAC}（差し込みも可） \\
ステーションの設定 & 通常のコントロール設定 & \textbf{タッチフリー用に設定を変える} \\
動作時間 & 実測から計算（Step 8） & \textbf{12 時間}が標準。
  タッチフリーのパンチでは動作時間が更新されないため \\
電池の減り & 通常 & \textbf{約 10 倍}。撤収後はすぐスタンバイに戻す \\
当日朝の追加作業 & なし & \textbf{SIAC のバッテリーテスト}。
  会場とスタートに \textbf{SIAC TEST} を設置 \\
競技者への周知 & 音と光の確認 & 電源の ON/OFF，
  反応しないときは差し込みでよいこと \\
コース設計 & 制約なし & \textbf{フィニッシュのそばを通らない}
  （通るとその場で電源が切れる） \\
\bottomrule
\end{tabular}
\end{center}

\begin{Note}[混ぜて使える]
同じ大会で\textbf{「選手権クラスは SIAC，一般クラスは SI カード」}という運用ができます。
タッチフリー設定のステーションは，通常の差し込みパンチも受け付けるからです。
SI カードのパンチがタッチフリーの妨げになることもありません。

逆にいうと，\textbf{SIAC を使う大会では，全ステーションをタッチフリー設定にしておけば，
どちらのカードの人も走れます}。
\end{Note}

\begin{Note}[押さえておきたい 3 点]
\begin{itemize}
\item \tEMIT\ ステーションの時計合わせは不要。そのかわり\textbf{スタートユニットの故障}が
      全員のタイムに効く。
\item \tSI\ ステーションの時計合わせは当日の朝に行う。とくにスタートとフィニッシュ。
\item \tSIAC\ フィニッシュをパンチすると SIAC の電源が切れる。
      そのため，コースがフィニッシュのそばを通らないように設計する。
\end{itemize}
\end{Note}

\begin{Note}[どのくらいの精度で計れていればよいか]
1 秒単位で成績を出すには，0.5 秒以上の精度で計れている必要があります。
スプリントなどで 0.1 秒単位を出すなら 0.05 秒の精度が要り，
これは \tSIAC\ タッチフリーや光電管を使わないと届きません。
通常の大会では，
\textbf{スタート地区の時計・チャイマー・\tSI\ フィニッシュのステーション}が
そろっていれば 1 秒単位で問題のない精度になります。
逆に，この 3 つのどれかがずれていると，競技者から見て分かる差が出ます（Step 24）。
\end{Note}

\section{当日のタイムライン（全体像）}
\tALL

第 3 部の Step は，このタイムラインに沿って並んでいます。

\begin{Fig}
\begin{tikzpicture}[x=1mm,y=1mm]
  \draw[line width=1pt,rulec] (0,0)--(160,0);
  \foreach \x/\t/\d in {
     4/開場前/{ネットワーク接続\\時計合わせ\\ポ確カード読み取り},
     40/開場〜受付/{当日申込入力\\カード変更\\当日版リスト配達},
     78/競技中/{カード読み取り\\ペナチェック\\印刷・速報},
    116/スタート閉鎖後/{欠席（DNS）入力\\表彰対象の報告},
    146/閉鎖・撤収/{未帰還者確認\\レンタル回収}}{
    \fill[boxln] (\x,0) circle (1.1);
    \node[font=\sffamily\bfseries\scriptsize,above=1.5mm of {(\x,0)},anchor=south] {\t};
    \node[font=\sffamily\tiny,below=2mm of {(\x,0)},anchor=north,align=center] {\d};
  }
\end{tikzpicture}
\figcap{大会当日に計センがやることの流れ}
\end{Fig}

\begin{Fail}
\textbf{起きること：}「開場前」の作業を，開場後に始めることになる。\\
\textbf{対処：}会場が開くと，参加者対応・当日申込・受付からの届出が同時に来ます。
落ち着いて作業できるのは開場前だけなので，
ネットワーク・時計合わせ・ポ確の 3 つはここで終わらせる前提で朝の時間を組みます。
間に合わないときは時計合わせを先にします。あとから直せないのはこれだけです。
\end{Fail}
